\section{Mathematical model}
\label{sec:model}

\subsection{Data}
\label{sec:data}

Let $V$ be the set of routers and $A$ the set of directed links in the nominal network. Every link
$a=(u,v)\in A$ has a strictly positive administrative metric $\omega_a$ and a capacity $c_a>0$. Let
$D$ be the set of traffic demands. A demand $d\in D$ is specified by its source $s_d\in V$,
destination $\tau_d\in V$, and nonnegative volume $v_{dt}$ in period~$t$. The experiments use either
$T=\{0\}$ or $T=\{0,1\}$. For each period, $q_t\subseteq A$ denotes the links removed by the
intervention and
\[
A_t=A\setminus q_t,\qquad G_t=(V,A_t).
\]
In the supplied instances $q_0=\emptyset$. The scenario also provides the maximum number of segments
$M=\mathrm{maxSeg}=6$ and, for the transition from period~0 to period~1, a reconfiguration
budget~$\kappa$.

These parameters are read from the three files required by the project specification: the network
file supplies $(V,A,\omega,c)$, the traffic-matrix file supplies $(D,s,\tau,v)$, and the scenario
file supplies $(q,M,\kappa)$. The JSON representation is therefore only the input format; the
optimization model below is defined in terms of the mathematical sets and parameters above. Node and
link identifiers are mapped explicitly rather than assumed to be consecutive.

Only a reachable ordered pair $(i,j)$ with $i\neq j$ can be used as a segment in period~$t$; let
$S_t$ denote the set of all such pairs. Between the endpoints of a segment, traffic is forwarded on
shortest paths in $G_t$. Let $\ell^t_{ij}$ be the shortest-path distance from $i$ to $j$ and let
$\sigma^t_{ij}$ be the number of shortest directed paths. For an available link $a=(u,v)\in A_t$, the
split coefficient $r^t_{ija}$ is the fraction of the traffic of segment $(i,j)$ carried by~$a$. The
course note~\citep{courser2026} gives
\begin{equation}
r^t_{ija}=
\begin{cases}
\dfrac{\sigma^t_{iu}\,\sigma^t_{vj}}{\sigma^t_{ij}}, & \text{if } \ell^t_{iu}+\omega_a+\ell^t_{vj}=\ell^t_{ij},\\[0.6em]
0, & \text{otherwise.}
\end{cases}
\label{eq:r}
\end{equation}
All distances and path counts, and hence all $r^t_{ija}$, are recomputed after an intervention
because deleting a link can change both reachability and the shortest-path DAG. The implementation
runs Dijkstra's algorithm from every router and stores path counts in logarithmic form to avoid
numerical overflow.

Equation~\eqref{eq:r} assigns equal probability to complete shortest paths and is used throughout the
experiments. It is not identical to a router-local ECMP recursion, which can assign different
probabilities to complete shortest paths on an asymmetric DAG. Consequently, the numerical loads in
this report are internally consistent with the supplied course note but may differ from an evaluator
that applies router-local ECMP.

\subsection{Decision variables}

For every $d\in D$, $t\in T$, and $(i,j)\in S_t$, define
\[
x^{dt}_{ij}=
\begin{cases}
1, & \text{if the segment path of demand } d \text{ contains segment } (i,j) \text{ in period } t,\\
0, & \text{otherwise.}
\end{cases}
\]
Thus, the nonzero $x$ variables describe the ordered waypoint sequence selected for each demand and
period; the physical links used between two consecutive waypoints are induced by the split
coefficients rather than selected directly.

The continuous variable $L\ge 0$ is an upper bound on every available-link utilization in every
modeled period. No separate load variable is required because each load is a linear expression
in~$x$. For the two-period model, an auxiliary continuous variable $y^d_{ij}\ge 0$ represents the
absolute change of segment $(i,j)$ for demand~$d$ between periods.

\subsection{Objective}
\label{sec:objective}

For a link $a\in A_t$, define its carried flow and utilization as
\begin{equation}
F_{at}=\sum_{d\in D}\sum_{(i,j)\in S_t} r^t_{ija}\,v_{dt}\,x^{dt}_{ij},
\qquad
\lambda_{at}=\frac{F_{at}}{c_a}.
\label{eq:load-def}
\end{equation}
The implemented objective is the scalar maximum link utilization (MLU):
\begin{equation}
\min L,\qquad L=\max_{t\in T,\,a\in A_t}\lambda_{at}.
\label{eq:mlu}
\end{equation}
This classical intra-domain traffic-engineering criterion suppresses bottlenecks, distributes traffic
over the available capacity, and preserves residual capacity for growth, bursts, and
failures~\citep{fortz2002te,texcp2005}. Because queuing delay rises nonlinearly as utilization
approaches capacity, controlling the peak is more protective of worst-case performance than minimizing
average load alone~\citep{fortz2002te}.

The challenge objective~\citep{orange2026} is stronger: after minimizing the largest element of the
load vector, it minimizes the second largest, then the third largest, and so on. Objective
\eqref{eq:mlu} is exactly the first stage of that lexicographic objective, but it does not break ties
between solutions having the same MLU. Accordingly, every lower bound, gap, and ``MLU-optimal'' label
in this report refers only to the scalar problem~\eqref{eq:mlu}.

\subsection{Constraints}

\paragraph{Flow conservation.}
For each demand and period, the selected segments must form a unit flow from its source to its
destination:
\begin{equation}
\sum_{j:(i,j)\in S_t} x^{dt}_{ij}-\sum_{j:(j,i)\in S_t} x^{dt}_{ji}=
\begin{cases}
1, & i=s_d,\\
-1, & i=\tau_d,\\
0, & \text{otherwise,}
\end{cases}
\qquad \forall d\in D,\ t\in T,\ i\in V.
\label{eq:flow}
\end{equation}
The formulation also adds the valid degree inequalities
\begin{equation}
\sum_{j:(i,j)\in S_t} x^{dt}_{ij}\le \mathbf{1}\{i\neq\tau_d\},
\qquad
\sum_{j:(j,i)\in S_t} x^{dt}_{ji}\le \mathbf{1}\{i\neq s_d\},
\qquad \forall d,t,i.
\label{eq:degree}
\end{equation}
They forbid branching, forbid a segment entering the source or leaving the destination, and make the
source--destination component a simple path. A disconnected directed cycle is still formally
possible, but deleting it preserves conservation and cannot increase $L$, since every traffic
coefficient is nonnegative.

\paragraph{Segment limit.}
The complete segment path, including its first and last segments, is limited by
\begin{equation}
\sum_{(i,j)\in S_t} x^{dt}_{ij}\le M,\qquad \forall d\in D,\ t\in T.
\label{eq:segments}
\end{equation}
This is the bound on the segment list carried by the packet, not a bound on the number of physical
links traversed inside the shortest-path forwarding graphs.

\paragraph{Link utilization.}
Substituting~\eqref{eq:load-def} into $\lambda_{at}\le L$ gives the linear constraints
\begin{equation}
\sum_{d\in D}\sum_{(i,j)\in S_t} r^t_{ija}\,v_{dt}\,x^{dt}_{ij}\le c_a L,
\qquad \forall t\in T,\ a\in A_t.
\label{eq:utilization}
\end{equation}
The left-hand side sums the fractions of all selected segments that traverse link~$a$. Dividing by
$c_a$ shows directly that~\eqref{eq:utilization} bounds every link utilization by $L$; because $L$ is
minimized, at least one such inequality is tight at an optimum. Loads above one remain allowed and
indicate congestion.

\paragraph{Reconfiguration budget.}
Let $S=S_0\cup S_1$ and define $\bar x^{dt}_{ij}=x^{dt}_{ij}$ when $(i,j)\in S_t$, and
$\bar x^{dt}_{ij}=0$ otherwise. For each demand and possible segment, the absolute difference is
linearized as
\begin{equation}
y^d_{ij}\ge \bar x^{d,1}_{ij}-\bar x^{d,0}_{ij},
\qquad
y^d_{ij}\ge \bar x^{d,0}_{ij}-\bar x^{d,1}_{ij},
\qquad \forall d\in D,\ (i,j)\in S.
\label{eq:change}
\end{equation}
The total number of segment-usage changes is bounded by
\begin{equation}
\sum_{d\in D}\sum_{(i,j)\in S} y^d_{ij}\le\kappa.
\label{eq:budget}
\end{equation}
The two inequalities force $y^d_{ij}\ge\lvert\bar x^{d,1}_{ij}-\bar x^{d,0}_{ij}\rvert$; any
unnecessary excess only consumes budget and can be reduced without affecting another constraint.
Hence~\eqref{eq:budget} bounds the symmetric difference between the two segment sets exactly as
required by the project definition. It is omitted in the period-0 experiment and activated only in the
joint two-period model.
